\chapter*{Abstract}
%\addcontentsline{toc}{chapter}{Abstract}
\label{chap:abstract}
\pagestyle{plain}
This report presents the Genetic Algorithm (GA) milestone of our project on cooperative spacecraft
fleets that photograph the surface of an irregular asteroid, and compares it with the Simulated
Annealing (SA) of Milestone~3. The optimiser chooses how many spacecraft to send (one on/off switch
per slot), the circular orbit of each and the mission duration; the fitness rewards the coverage of
the sunlit surface and penalises fleet cost, energy and mission time, under orbit-radius, separation
and battery constraints. Every candidate is scored by a true-scale mission simulation with Kepler
orbits, a spinning non-convex shape model, ray-cast occlusion, sunlight, a camera resolution limit
and batteries.
The GA is written from scratch on the same Algorithm $\rightarrow$ Feasibility $\rightarrow$ Engine
pipeline. Its chromosome is the whole mission, one five-gene slot per spacecraft. Problem-specific
operators align the slots of two parents before crossover (a fleet can be stored in $N_{\max}!$
orders), swap whole spacecraft or blend matched orbits with BLX-$\alpha$ along the shorter arc for the
angles, start from a population stratified over the fleet size, and reuse SA's move as mutation.
A parameter study of 23 settings confirmed every default (population 40 or 60, 5\,\% elites, crossover
fraction 0.8, binary tournament, blend probability 0.3) and showed that both crossover and mutation are
needed and that midpoint averaging is the worst choice ($p=0.004$).
The GA finds the known optima of benchmark functions in every run, reproduces the 98\,\% literature
coverage of Eros, and beats random search in every case. On nine case studies (26 or 51 genes,
2000 evaluations, five seeds) SA is slightly better (mean fitness 0.266 against 0.273, significant in
two cases) and converges about 1.7 times faster, while the GA is on a par in seven cases, varies half
as much between runs, and sends slightly larger fleets with equal or higher coverage. On harder tests
the ranking turns: the GA is better on all three short Kleopatra missions, keeps finding small fleets
with 101 genes where SA gets stuck at 14--15 spacecraft, and with 5000 evaluations closes the gap and
becomes significantly better on the hardest Eros mission ($p=0.048$).
